\cleardoublepage
\vspace{2cm}
\begin{center}
{\fontsize{14}{16.8}\selectfont\textbf{DAFTAR SINGKATAN}}\\[1em]
\end{center}
\label{chap:singkatan}
\addcontentsline{toc}{chapter}{DAFTAR SINGKATAN}
\setcounter{chapter}{0}
\setcounter{section}{0}
\vspace{2em}

% Singkatan berikut didaftar berdasarkan kemunculan aktual pada Bab 1--5.
% Urutan ditampilkan alfabetis. Kolom Bab menunjuk lokasi pemakaian pertama.

\begin{longtable}{@{} l >{\raggedright\arraybackslash}p{8cm} c @{}}
\multicolumn{1}{l}{\textbf{Singkatan}} & \multicolumn{1}{l}{\textbf{Deskripsi}} & \multicolumn{1}{l}{\parbox[c]{2.5cm}{\raggedright \textbf{Pemakaian pertama kali}}} \\
\hline

ACID & \textit{Atomicity, Consistency, Isolation, Durability}, sifat transaksi basis data. & Bab~\ref{chap:studi-literatur} \\
AGPL & \textit{GNU Affero General Public License}, lisensi \textit{copyleft} kuat yang juga mencakup penggunaan melalui jaringan. & Bab~\ref{chap:studi-literatur} \\
ANN & \textit{Approximate Nearest Neighbor}, kelas algoritma pencarian tetangga terdekat secara aproksimasi pada ruang vektor. & Bab~\ref{chap:studi-literatur} \\
API & \textit{Application Programming Interface}. & Bab~\ref{chap:pendahuluan} \\
APISIX & Apache APISIX, \textit{API gateway} berbasis Nginx dan etcd. & Bab~\ref{chap:studi-literatur} \\
ASF & \textit{Apache Software Foundation}, yayasan nirlaba penaung proyek \textit{open source} Apache. & Bab~\ref{chap:studi-literatur} \\
BSD & \textit{Berkeley Software Distribution}, keluarga lisensi permisif perangkat lunak. & Bab~\ref{chap:studi-literatur} \\
BSL & \textit{Business Source License}, lisensi \textit{source-available} dengan pembatasan pemakaian komersial. & Bab~\ref{chap:studi-literatur} \\
BUSL & Penulisan SPDX untuk \textit{Business Source License} versi 1.1. & Bab~\ref{chap:studi-literatur} \\
CD & \textit{Continuous Delivery} atau \textit{Continuous Deployment}. & Bab~\ref{chap:pendahuluan} \\
CDC & \textit{Change Data Capture}, teknik penangkapan perubahan data pada basis data sumber. & Bab~\ref{chap:studi-literatur} \\
CI & \textit{Continuous Integration}. & Bab~\ref{chap:studi-literatur} \\
CNCF & \textit{Cloud Native Computing Foundation}, yayasan penaung proyek \textit{cloud native} di bawah Linux Foundation. & Bab~\ref{chap:studi-literatur} \\
CNPG & \textit{CloudNativePG}, \textit{operator} Kubernetes untuk PostgreSQL. & Bab~\ref{chap:implementasi} \\
CPU & \textit{Central Processing Unit}. & Bab~\ref{chap:studi-literatur} \\
CT & \textit{Continuous Training}, pelatihan model secara berkelanjutan dengan pemicu \textit{drift} atau jadwal. & Bab~\ref{chap:studi-literatur} \\
DSRM & \textit{Design Science Research Methodology}, kerangka metodologi penelitian rekayasa artefak. & Bab~\ref{chap:pendahuluan} \\
eBPF & \textit{extended Berkeley Packet Filter}, mekanisme menjalankan program tervalidasi di dalam kernel Linux. & Bab~\ref{chap:studi-literatur} \\
ESO & \textit{External Secrets Operator}, \textit{operator} Kubernetes penyinkron rahasia dari pengelola rahasia eksternal. & Bab~\ref{chap:studi-literatur} \\
GPU & \textit{Graphics Processing Unit}. & Bab~\ref{chap:studi-literatur} \\
HNSW & \textit{Hierarchical Navigable Small World}, struktur graf untuk ANN. & Bab~\ref{chap:studi-literatur} \\
HPA & \textit{Horizontal Pod Autoscaler}, kontrol penskalaan jumlah \textit{pod} pada Kubernetes. & Bab~\ref{chap:pendahuluan} \\
HTML & \textit{HyperText Markup Language}. & Bab~\ref{chap:perancangan} \\
JSON & \textit{JavaScript Object Notation}. & Bab~\ref{chap:studi-literatur} \\
KEDA & \textit{Kubernetes Event-driven Autoscaling}, \textit{autoscaler} berbasis sinyal eksternal. & Bab~\ref{chap:pendahuluan} \\
KES & \textit{Key Encryption Service}, layanan kunci enkripsi sisi server untuk MinIO. & Bab~\ref{chap:studi-literatur} \\
KS & \textit{Kolmogorov-Smirnov test}, uji statistik kesetaraan dua distribusi. & Bab~\ref{chap:studi-literatur} \\
LF & \textit{Linux Foundation}, yayasan nirlaba penaung ekosistem proyek \textit{open source}. & Bab~\ref{chap:studi-literatur} \\
MLOps & \textit{Machine Learning Operations}, disiplin praktik operasional pengembangan dan penyajian model \textit{machine learning}. & Bab~\ref{chap:pendahuluan} \\
MPL & \textit{Mozilla Public License}, lisensi \textit{copyleft} lemah berbasis berkas. & Bab~\ref{chap:studi-literatur} \\
OIDC & \textit{OpenID Connect}, protokol identitas berbasis OAuth~2.0. & Bab~\ref{chap:studi-literatur} \\
OLAP & \textit{Online Analytical Processing}, pemrosesan kueri analitis pada data bervolume besar. & Bab~\ref{chap:studi-literatur} \\
ONNX & \textit{Open Neural Network Exchange}, format pertukaran model lintas \textit{framework}. & Bab~\ref{chap:perancangan} \\
PSI & \textit{Population Stability Index}, metrik pergeseran distribusi fitur antar dua periode. & Bab~\ref{chap:studi-literatur} \\
REST & \textit{Representational State Transfer}. & Bab~\ref{chap:studi-literatur} \\
SDK & \textit{Software Development Kit}. & Bab~\ref{chap:studi-literatur} \\
SLO & \textit{Service Level Objective}. & Bab~\ref{chap:studi-literatur} \\
SSO & \textit{Single Sign-On}, autentikasi tunggal untuk banyak aplikasi. & Bab~\ref{chap:studi-literatur} \\
TLS & \textit{Transport Layer Security}. & Bab~\ref{chap:pendahuluan} \\
VPA & \textit{Vertical Pod Autoscaler}, kontrol penyetelan \textit{request}/\textit{limit} sumber daya \textit{pod}. & Bab~\ref{chap:pendahuluan} \\

\end{longtable}
